
\documentclass{article}
\usepackage{colortbl}
\usepackage{makecell}
\usepackage{multirow}
\usepackage{supertabular}

\begin{document}

\newcounter{utterance}

\centering \large Interaction Transcript for game `cladder', experiment `full\_v1.5\_default', episode 4386 with selfrefine{-}Qwen3.5{-}9B.
\vspace{24pt}

{ \footnotesize  \setcounter{utterance}{1}
\setlength{\tabcolsep}{0pt}
\begin{supertabular}{c@{$\;$}|p{.15\linewidth}@{}p{.15\linewidth}p{.15\linewidth}p{.15\linewidth}p{.15\linewidth}p{.15\linewidth}}
    \# & \multicolumn{2}{c}{Player} && \multicolumn{2}{c}{Game Master} \\
    \hline

    \theutterance \stepcounter{utterance}  
    & & & \multicolumn{4}{p{0.6\linewidth}}{
        \cellcolor[rgb]{0.9,0.9,0.9}{
            \makecell[{{p{\linewidth}}}]{
                \texttt{\tiny{[P1$\langle$GM]}}
                \texttt{You are an expert in causal inference. The following question is not a typical commonsense query, but rather a meticulously designed question created by a professor specializing in causal inference, intended to assess the students' mastery of the course content.} \\
\texttt{Imagine a self{-}contained, hypothetical world with only the following conditions, and without any unmentioned factors or causal relationships: Gender has a direct effect on smoking and peanut allergy. Smoking has a direct effect on tar deposit. Tar deposit has a direct effect on peanut allergy. Gender is unobserved. For nonsmokers and with no tar deposit, the probability of being allergic to peanuts is 35\%. For nonsmokers and with high tar deposit, the probability of being allergic to peanuts is 63\%. For smokers and with no tar deposit, the probability of being allergic to peanuts is 25\%. For smokers and with high tar deposit, the probability of being allergic to peanuts is 49\%. For nonsmokers, the probability of high tar deposit is 22\%. For smokers, the probability of high tar deposit is 73\%. For smokers, would it be more likely to see being allergic to peanuts if the person had been a nonsmoker?} \\
\texttt{Start your answer with 'Yes' or 'No', followed by additional reasoning or evidence to support your explanation.} \\
            }
        }
    }
     \\ \\

    \theutterance \stepcounter{utterance}  
    & & \multicolumn{4}{p{0.6\linewidth}}{
        \cellcolor[rgb]{0.95,0.95,0.95}{
            \makecell[{{p{\linewidth}}}]{
                \texttt{\tiny{[P1$\leftrightarrow$P]}}
                \texttt{{[}DRAFT{]}} \\
\texttt{No} \\
\\ 
\texttt{**Reasoning:**} \\
\\ 
\texttt{To determine if a smoker would be more likely to be allergic to peanuts if they had been a nonsmoker, we must calculate the **Counterfactual Probability** \$P(\textbackslash{}text\{Allergy\} \textbackslash{}mid \textbackslash{}text\{Smoker\}, \textbackslash{}text\{Nonsmoker\})\$. This requires estimating the probability of having a peanut allergy for a smoker under the hypothetical condition that they do not smoke.} \\
\\ 
\texttt{In this causal structure, the effect of smoking on peanut allergy is entirely mediated through **tar deposit**. There is no direct path from Smoking to Allergy that bypasses Tar Deposit. Therefore, the counterfactual probability for a smoker who did not smoke is equal to the probability of having a peanut allergy for a **nonsmoker** with the same distribution of tar deposits.} \\
\\ 
\texttt{However, the prompt provides specific conditional probabilities for allergies based on smoking status and tar deposit levels, but it does not explicitly state that the distribution of tar deposits is independent of gender or that the "nonsmoker" group in the counterfactual scenario has the same tar deposit distribution as the actual smoker.} \\
\\ 
\texttt{Let's analyze the causal mechanism strictly based on the provided numbers:} \\
\texttt{1.  **Actual Scenario (Smoker):**} \\
\texttt{    *   \$P(\textbackslash{}text\{High Tar\} \textbackslash{}mid \textbackslash{}text\{Smoker\}) = 0.73\$} \\
\texttt{    *   \$P(\textbackslash{}text\{Low Tar\} \textbackslash{}mid \textbackslash{}text\{Smoker\}) = 0.27\$} \\
\texttt{    *   \$P(\textbackslash{}text\{Allergy\} \textbackslash{}mid \textbackslash{}text\{Smoker\}, \textbackslash{}text\{High Tar\}) = 0.49\$} \\
\texttt{    *   \$P(\textbackslash{}text\{Allergy\} \textbackslash{}mid \textbackslash{}text\{Smoker\}, \textbackslash{}text\{Low Tar\}) = 0.25\$} \\
\texttt{    *   Actual \$P(\textbackslash{}text\{Allergy\} \textbackslash{}mid \textbackslash{}text\{Smoker\}) = (0.73 \textbackslash{}times 0.49) + (0.27 \textbackslash{}times 0.25) = 0.3577 + 0.0675 = 0.4252\$ (42.52\%).} \\
\\ 
\texttt{2.  **Counterfactual Scenario (Smoker had been a Nonsmoker):**} \\
\texttt{    *   If a person stops smoking (becomes a nonsmoker), their probability of having high tar deposit changes. The prompt states: "For nonsmokers, the probability of high tar deposit is 22\%."} \\
\texttt{    *   Crucially, in a causal world where "Smoking has a direct effect on tar deposit," the counterfactual state of "Nonsmoker" implies the probability of high tar deposit reverts to the nonsmoker baseline (\$0.22\$), not the smoker baseline (\$0.73\$).} \\
\texttt{    *   We must use the allergy probabilities} \\
            }
        }
    }
    & \\ \\

    \theutterance \stepcounter{utterance}  
    & & \multicolumn{4}{p{0.6\linewidth}}{
        \cellcolor[rgb]{0.95,0.95,0.95}{
            \makecell[{{p{\linewidth}}}]{
                \texttt{\tiny{[P1$\leftrightarrow$P]}}
                \texttt{{[}SELF{-}FEEDBACK{]}} \\
\texttt{FORMAT: problem} \\
\texttt{RULES: ok} \\
\texttt{STRATEGY: problem} \\
\texttt{REASON: The answer fails to complete the calculation required to determine the counterfactual probability and incorrectly concludes the answer is "No" without showing that the nonsmoker baseline probability (35\%) is lower than the actual smoker probability (\textasciitilde{}42.5\%).} \\
            }
        }
    }
    & \\ \\

    \theutterance \stepcounter{utterance}  
    & \multicolumn{4}{p{0.6\linewidth}}{
        \cellcolor[rgb]{0.561,0.737,0.561}{
            \makecell[{{p{\linewidth}}}]{
                \texttt{\tiny{[P1$\rangle$GM]}}
                \texttt{problem} \\
\texttt{No} \\
\\ 
\texttt{**Reasoning:**} \\
\\ 
\texttt{To determine if a smoker would be more likely to be allergic to peanuts if they had been a nonsmoker, we must compare the actual probability of allergy for a smoker against the counterfactual probability of allergy for that same individual had they been a nonsmoker.} \\
\\ 
\texttt{1.  **Calculate the Actual Probability (\$P(\textbackslash{}text\{Allergy\} \textbackslash{}mid \textbackslash{}text\{Smoker\})\$):**} \\
\texttt{    Using the law of total probability with the provided conditional probabilities for tar deposits among smokers:} \\
\texttt{    *   \$P(\textbackslash{}text\{High Tar\} \textbackslash{}mid \textbackslash{}text\{Smoker\}) = 0.73\$} \\
\texttt{    *   \$P(\textbackslash{}text\{Low Tar\} \textbackslash{}mid \textbackslash{}text\{Smoker\}) = 1 {-} 0.73 = 0.27\$} \\
\texttt{    *   \$P(\textbackslash{}text\{Allergy\} \textbackslash{}mid \textbackslash{}text\{Smoker\}, \textbackslash{}text\{High Tar\}) = 0.49\$} \\
\texttt{    *   \$P(\textbackslash{}text\{Allergy\} \textbackslash{}mid \textbackslash{}text\{Smoker\}, \textbackslash{}text\{Low Tar\}) = 0.25\$} \\
\\ 
\texttt{    \$P(\textbackslash{}text\{Allergy\} \textbackslash{}mid \textbackslash{}text\{Smoker\}) = (0.73 \textbackslash{}times 0.49) + (0.27 \textbackslash{}times 0.25) = 0.3577 + 0.0675 = 0.4252\$ (42.52\%).} \\
\\ 
\texttt{2.  **Calculate the Counterfactual Probability (\$P(\textbackslash{}text\{Allergy\} \textbackslash{}mid \textbackslash{}text\{Smoker\}, \textbackslash{}text\{Nonsmoker\})\$):**} \\
\texttt{    In this causal structure, the effect of smoking on allergy is fully mediated by tar deposit. If a smoker had been a nonsmoker, their probability of having high tar deposit would revert to the baseline for nonsmokers (\$0.22\$), and their allergy probability would follow the nonsmoker conditional rates.} \\
\texttt{    *   \$P(\textbackslash{}text\{High Tar\} \textbackslash{}mid \textbackslash{}text\{Nonsmoker\}) = 0.22\$} \\
\texttt{    *   \$P(\textbackslash{}text\{Low Tar\} \textbackslash{}mid \textbackslash{}text\{Nonsmoker\}) = 1 {-} 0.22 = 0.78\$} \\
\texttt{    *   \$P(\textbackslash{}text\{Allergy\} \textbackslash{}mid \textbackslash{}text\{Nonsmoker\}, \textbackslash{}text\{High Tar\}) = 0.63\$} \\
\texttt{    *   \$P(\textbackslash{}text\{Allergy\} \textbackslash{}mid \textbackslash{}text\{Nonsmoker\}, \textbackslash{}text\{Low Tar\}) = 0.35\$} \\
\\ 
\texttt{    \$P(\textbackslash{}text\{Allergy\} \textbackslash{}mid \textbackslash{}text\{Nonsmoker\}) = (0.22 \textbackslash{}times 0.63) + (0.78 \textbackslash{}times} \\
            }
        }
    }
    & & \\ \\

    \theutterance \stepcounter{utterance}  
    & & & \multicolumn{2}{p{0.3\linewidth}}{
        \cellcolor[rgb]{0.9,0.9,0.9}{
            \makecell[{{p{\linewidth}}}]{
                \texttt{\tiny{[GM$|$GM]}}
                \texttt{ParseError: The response doesn't start with one of the required word {[}'yes', 'no'{]}, but problem} \\
            }
        }
    }
    & & \\ \\

    \theutterance \stepcounter{utterance}  
    & & & \multicolumn{2}{p{0.3\linewidth}}{
        \cellcolor[rgb]{0.9,0.9,0.9}{
            \makecell[{{p{\linewidth}}}]{
                \texttt{\tiny{[GM$|$GM]}}
                \texttt{game\_result = ABORT} \\
            }
        }
    }
    & & \\ \\

\end{supertabular}
}

\end{document}
